\section*{Chapter \ref{ch:rc:credit-and-debit-valuation-adjustments} --- Credit and Debit Valuation Adjustments}

\begin{solution}{exo:rc:credit-and-debit-valuation-adjustments:1}
$0.6\times2\,000\,000\times(1-e^{-0.02\times5}) = \text{USD}~114\,195$.
\end{solution}

\begin{solution}{exo:rc:credit-and-debit-valuation-adjustments:2}
$6.44\text{ million}\times0.0175\times7.45 = \text{USD}~840\,420$, against the exact 891\,342: the
approximation uses an undiscounted, time-averaged exposure and the ten-year spread for every
date, and it works because the spread already prices $\lambda(1-R)$.
\end{solution}

\begin{solution}{exo:rc:credit-and-debit-valuation-adjustments:3}
The \omterm{def:rc:counterparty-exposure:netting}{netting set} drifts in the counterparty's favour, so the bank expects to owe more than it
is owed: its DVA exceeds its CVA even though the counterparty is riskier. The counterparty sees
the mirror image: its CVA on the bank is USD~896\,328 on first-to-default, its DVA 832\,416, a
net cost of 63\,912 that is the bank's net benefit.
\end{solution}

\begin{solution}{exo:rc:credit-and-debit-valuation-adjustments:4}
CVA is linear in the \omterm{def:rc:counterparty-exposure:profiles}{expected positive exposure}, and the positive part of a sum is at most
the sum of the positive parts: netting lets the two trades' values offset on each path.
\end{solution}

\begin{solution}{exo:rc:credit-and-debit-valuation-adjustments:5}
$4\,254/(7.45\times10^{-4}) \approx \text{USD}~5.7$ million of ten-year protection. It leaves the
changes of the exposure with rates and FX (hedged separately), \omterm{def:rc:rates-risk:crossgamma}{cross-gamma} between credit and
exposure, \omterm{def:rc:counterparty-exposure:wwr}{wrong-way risk}, the basis between the counterparty's curve and any proxy, and
jump-to-default: protection sized on \omterm{def:rc:reduced-form-credit:cs01}{CS01} does not pay the actual loss at default.
\end{solution}

\begin{solution}{exo:rc:credit-and-debit-valuation-adjustments:6}
Under the CSA both adjustments are driven by the value's moves over the margin period of
risk, which are symmetric; without it, the \omterm{def:rc:counterparty-exposure:netting}{netting set}'s drift in the counterparty's favour
made the negative exposure large. The CSA removes more of the larger side.
\end{solution}

\begin{solution}{exo:rc:credit-and-debit-valuation-adjustments:7}
USD~766\,532, 1\,108\,459 and 1\,521\,524. With a very strong dependence, paths where the euro falls
early default almost surely in the first years, when the exposure is still small, and the
average survival must be kept, so the defaults are taken from the later, larger exposures: the
CVA turns down after $b\approx7$.
\end{solution}

\begin{solution}{exo:rc:credit-and-debit-valuation-adjustments:8}
A DVA gain says the bank's own credit worsened: the market now thinks it likelier to default
on what it owes. It is value to creditors realised only by defaulting, not a sign of a better
business; that is why it is deducted from capital and usually excluded from traders' P\&L.
\end{solution}

\begin{solution}{pb:rc:credit-and-debit-valuation-adjustments:1}
\textbf{1.} USD~1\,055\,437 before, 1\,513\,008 after.
\textbf{2.} USD~457\,571, a gain in the fair value of derivatives through profit and loss.
\textbf{3.} The liabilities (the negative exposure) are now less likely to be paid in full, so
their fair value falls: the bank gains as a debtor.
\textbf{4.} About USD~229 million.
\textbf{5.} The bank's counterparties: their CVA on the bank rises by the same amount.
\textbf{6.} It deducts the DVA from common equity tier~1: the gain does not count as capital.
\textbf{7.} From 1 January 2018, after a phase-in from 20\% in 2014.
\textbf{8.} CVA is a loss that a counterparty's default would realise; DVA is a gain realised only
by the bank's own default, when capital is needed and absent.
\textbf{9.} The ratio would rise as the bank weakened, the opposite of what capital should do.
\textbf{10.} As a non-recurring item to strip out: most analysts report results before own-credit
adjustments.
\textbf{11.} A default swap on oneself would pay only when the bank has defaulted and cannot
collect; no one would sell it at a useful price.
\textbf{12.} Selling protection on peer banks, which gains when bank spreads tighten together;
it adds the peers' \omterm{def:rc:reduced-form-credit:jtd}{jump-to-default risk} and the basis between the bank's and the peers' spreads.
\textbf{13.} Most banks do not: the hedge adds real risk to protect an accounting number that
capital already excludes.
\textbf{14.} Under the zero-threshold CSA the DVA is USD~70\,490, so the gain from the same
widening is about 93\% smaller.
\textbf{15.} Its CVA on the bank rises by the same USD~457\,571, a loss in its accounts.
\textbf{16.} Value to the bank's creditors (they bear the shortfall), not to its shareholders.
\textbf{17.} Report it separately and exclude it from the desk's performance, or charge it to a
central function; traders should not be rewarded for the bank's credit deteriorating.
\textbf{18.} Each party's \omterm{def:rc:credit-and-debit-valuation-adjustments:bcva}{bilateral CVA} is the other's bilateral DVA, so a price that includes
both is symmetric; a unilateral price counts the same default twice or not at all.
\textbf{19.} Named result: \emph{the gain from getting riskier}: a 50-basis-point widening raises the
\omterm{def:rc:counterparty-exposure:netting}{netting set}'s DVA from USD~1\,055\,437 to 1\,513\,008, a gain of USD~457\,571.
\textbf{20.} Because it grows exactly when the bank's ability to absorb losses shrinks.
\end{solution}

\begin{solution}{iq:rc:credit-and-debit-valuation-adjustments:1}
$\mathrm{CVA} = (1-R)\sum_k\mathrm{DEE}(t_k)\,\mathrm{PD}(t_{k-1},t_k)$: \omterm{def:rc:reduced-form-credit:pd}{loss given default}, discounted
\omterm{def:rc:counterparty-exposure:profiles}{expected positive exposure} of the \omterm{def:rc:counterparty-exposure:netting}{netting set} at each date, and the risk-neutral probability of
the counterparty defaulting in each period from its default-swap curve; independence assumed.
\iqlookfor{the three terms and the independence assumption.}
\end{solution}

\begin{solution}{iq:rc:credit-and-debit-valuation-adjustments:2}
The counterparty's CVA on the bank: the value of the bank's option to default on what it owes.
Controversial because it rises when the bank's credit worsens, can be realised only by default,
cannot be hedged directly, and flatters earnings in bad times; regulators deduct it from capital.
\iqlookfor{the definition and the three objections.}
\end{solution}

\begin{solution}{iq:rc:credit-and-debit-valuation-adjustments:3}
Hedge the exposure drivers (rates delta and vega) with swaps and swaptions, the credit with
default swaps on the name or a proxy index sized on \omterm{def:rc:reduced-form-credit:cs01}{CS01}, and manage jump-to-default with
single-name protection where it trades; seek a CSA or break clauses; monitor \omterm{def:rc:rates-risk:crossgamma}{cross-gamma}.
\iqlookfor{market and credit hedges and their limits.}
\end{solution}

\begin{solution}{iq:rc:credit-and-debit-valuation-adjustments:4}
Simulate the counterparty's hazard with the market factors (a stochastic intensity correlated
with them, or a function of them) and compute pathwise default probabilities; or shift
exposures conditional on default; or use stress scenarios. Keep the average survival curve
calibrated to the market.
\iqlookfor{joint simulation and calibration of the marginal.}
\end{solution}

\begin{solution}{iq:rc:credit-and-debit-valuation-adjustments:5}
Most counterparties did not default, but their spreads widened sharply and exposures grew as
rates and FX moved: the CVA, marked to market, rose on every uncollateralised \omterm{def:rc:counterparty-exposure:netting}{netting set} at
once, while defaults hit only a few.
\iqlookfor{mark-to-market of credit versus realised defaults.}
\end{solution}

\begin{solution}{iq:rc:credit-and-debit-valuation-adjustments:6}
Risk-free close-out values the terminated trades without the survivor's credit; risky
close-out lets the replacement price include it. With a risky close-out, the survivor's DVA
remains in the \omterm{def:rc:credit-and-debit-valuation-adjustments:closeout}{close-out amount}, which changes the bilateral formulas and the recovery of the
defaulted party's estate.
\iqlookfor{the two conventions and their effect on DVA.}
\end{solution}
